% !TEX root = main.tex
% --- Abstract (single page, single column) ---
\begin{center}
  \begin{minipage}{0.75\textwidth}
    \begin{center}
      \textbf{\large Abstract}
    \end{center}
    \vspace{0.8em}
    \noindent Hand gesture recognition is a central problem in Human--Computer
    Interaction (HCI), where the goal is to classify sequences of three-dimensional
    coordinates $\mathbf{r}(t) = [x(t), y(t), z(t)]^{\top}$ recorded over discrete
    time. This project implements and compares four classifiers on the gesture
    datasets introduced by Huang et al.~\cite{Huang2019}, namely \emph{domain 1}
    (digits 0--9) and \emph{domain 4} (3D figures). Two baseline methods based
    on dynamic programming are first considered: Dynamic Time Warping (DTW)
    and an Edit-distance-based classifier, the latter relying on a preliminary
    vector-quantization step. Two state-of-the-art approaches are then
    investigated: the Three-Cents ($3\mathbb{c}$) recognizer of Caputo et al.~\cite{Caputo2018},
    a rotation-free 3D adaptation of the \$1 recognizer~\cite{Wobbrock2007}, and a Recurrent
    Neural Network (RNN). All methods are evaluated through leave-one-user-out
    (user-independent) and leave-one-gesture-sample-out (user-dependent)
    cross-validation procedures, and compared using hypothesis testing. This
    report details the theoretical framework of each method, the implementation
    strategy, and a rigorous statistical comparison of their accuracies.
  \end{minipage}
\end{center}
